% Source: https://artofproblemsolving.com/community/c3h2587410_expected_value_of_the_experimental_probability
% Stable identity: aops-topic-2587410-q01
\textbf{(a)} Write $X=I_1+I_2+\cdots+I_n$, where $I_k=1$ if $A$ occurs on the $k$th trial and $I_k=0$ otherwise. Then $E[I_k]=1\cdot p+0\cdot(1-p)=p$, and by linearity of expectation
\[E\left[\dfrac Xn\right]=\dfrac1n\left(E[I_1]+\cdots+E[I_n]\right)=\dfrac{np}{n}=\boxed{p}.\]
(Independence is not even needed for this.)

\textbf{(b)} No. For example, with one toss of a fair coin ($n=1$, $p=\dfrac12$), the experimental probability of heads is $0$ or $1$, never $\dfrac12$. The experimental probability is a random quantity whose \emph{average} is $p$; for large $n$ it is usually close to $p$ (law of large numbers). $\blacksquare$
